% Generated draft: requires mathematical and source review.
\documentclass[11pt]{article}
\usepackage[margin=25mm]{geometry}
\usepackage{fontspec}
\setmainfont{Latin Modern Roman}
\usepackage{amsmath,amssymb,mathtools,unicode-math}
\setmathfont{Latin Modern Math}
\usepackage{enumitem,needspace,xurl,hyperref}
\setlength{\parindent}{0pt}
\setlength{\parskip}{5pt}
\setlength{\emergencystretch}{3em}
\newcommand{\N}{\mathbb{N}}
\newcommand{\Z}{\mathbb{Z}}
\newcommand{\Q}{\mathbb{Q}}
\newcommand{\R}{\mathbb{R}}
\newcommand{\C}{\mathbb{C}}
\newcommand{\F}{\mathbb{F}}
\newcommand{\A}{\mathbb{A}}
\newcommand{\PP}{\mathbb{P}}
\newcommand{\E}{\mathbb{E}}
\newcommand{\eps}{\varepsilon}
\newcommand{\dd}{\,\mathrm{d}}
\begin{document}
\section*{UnsolvedMath ID: 1200035}
\textbf{Source problem code:} catalogue entry
\subsection{Definitions and mathematical statement}
Let \((X,\mu)\) be a standard Borel probability space and let \(E\subseteq X\times X\) be a symmetric irreflexive Borel relation of bounded degree, regarded as an undirected simple graph. It is a graphing when its counting-edge measure is invariant under reversing endpoints: for every nonnegative Borel function \(f\) on \(E\),
\[
\int_X\sum_{y:(x,y)\in E}f(x,y)\,\dd\mu(x)
=
\int_X\sum_{y:(x,y)\in E}f(y,x)\,\dd\mu(x).
\]
It is \(d\)-regular if every vertex has degree \(d\), up to the usual null-set convention, for a finite integer \(d\ge0\). Simplicity retains the source's exclusion of multiple edges.

In the measurable, modulo-null-sets interpretation supplied by the research summary, the question is whether there exist an invariant measurable conull set \(X_0\subseteq X\) and a measurable coloring of its undirected edges
\[
c:E(X_0)\longrightarrow\{1,\ldots,d+1\}
\]
such that distinct edges sharing a vertex have different colors. Invariance of \(X_0\) means it is a union of connected components. Thus each color class is a measurable matching, and all edges on \(X_0\) are colored.

The supplied summary reports an affirmative measurable \((d+1)\)-edge-coloring result, attributing it to Greb\'ik--Pikhurko (2020) and an extension by Greb\'ik. The literal question does not specify measurable versus Borel coloring, or whether coloring must hold on every vertex rather than a conull invariant set. Those distinctions must be checked before interpreting the reported theorem as an everywhere Borel conclusion; no such stronger conclusion is added here.
\subsection{Short English statement}
Can every simple \(d\)-regular graphing have its edges properly colored with \(d+1\) colors in a measurable way? The supplied summary reports yes; the null-set and Borel-versus-measurable conventions need review.
\subsection{Sources}
\begin{itemize}
\item Local UnsolvedMath registry, permanent ID \(1200035\), exact source code \texttt{catalogue entry}, ``Some Questions --- Question 35.''
\item Abert, \emph{Some questions} (2010), Question 35; \url{https://www.renyi.hu/~abert/questions.pdf}; manual transcription; recorded access 2026-07-29.
\item Supplied research summary: Greb\'ik--Pikhurko, \emph{Advances in Mathematics} \textbf{374} (2020), 107386, arXiv:1903.02657; extension attributed to Greb\'ik, arXiv:2303.16440. The precise scopes and relation of these cited versions have not been independently checked.
\end{itemize}
\end{document}
